\section{Proofs and derivations}\label{app:proofs}

Throughout, $(f_i,Y_i)_{i\le n}$ and $(f_j)_{j\le N}$ are i.i.d.\ from the same distribution, $g(f)=\E[Y\mid f]$, $v(f)=\Var(Y\mid f)$, and all second moments are finite. For $h$ with $\E h(f)^2<\infty$ let
$\hat\theta_h=\frac1n\sum_{i\le n}(Y_i-h(f_i))+\frac1N\sum_{j\le N}h(f_j)$.

\paragraph{Proof of Proposition~\ref{prop:main}.}
$\hat\theta_h$ is the sum of two independent sample means, hence unbiased for $\theta=\E Y$ with
$\Var(\hat\theta_h)=V(h):=\Var(Y-h)/n+\Var(h)/N$ exactly, and asymptotically normal as $n,N\to\infty$ by the Lindeberg--Feller CLT applied to each mean (with $n/N$ bounded).
Since $Y-g(f)$ is uncorrelated with every square-integrable function of $f$, $\Var(Y-h)=\E v+\Var(g-h)$. Write $h=cg+\delta$ with $c=N/(n+N)$ (so $1-c=n/(n+N)$). Then
$\Var(g-h)=(1-c)^2\Var g-2(1-c)\Cov(g,\delta)+\Var\delta$ and $\Var(h)=c^2\Var g+2c\Cov(g,\delta)+\Var\delta$, so
\[
V(h)=\frac{\E v}{n}+\Big[\frac{(1-c)^2}{n}+\frac{c^2}{N}\Big]\Var g+2\Big[\frac{c}{N}-\frac{1-c}{n}\Big]\Cov(g,\delta)+\Big(\frac1n+\frac1N\Big)\Var\delta .
\]
With $c=N/(n+N)$: $(1-c)^2/n+c^2/N=1/(n+N)$ and $c/N=(1-c)/n=1/(n+N)$, so the cross term vanishes and $V(h)=V_\star+(1/n+1/N)\Var(h-cg)$. \qed

\paragraph{Proof of Corollary~\ref{cor:lin}.}
For $h=\lambda f$, $\min_\lambda\Var(\lambda f-cg)=c^2\Var(g)(1-\rho^2)$ with $\rho=\mathrm{Corr}(g(f),f)$ (residual variance of the linear projection of $cg$ on $f$). By Proposition~\ref{prop:main}, $V_{\rm lin}-V_\star=(1/n+1/N)c^2\Var(g)(1-\rho^2)$, and $(1/n+1/N)c^2=\frac{n+N}{nN}\frac{N^2}{(n+N)^2}=\frac{N}{n(n+N)}$. Thus $V_{\rm lin}-V_\star=\frac{N}{n(n+N)}\Var(g)(1-\rho^2)$.
For the classical estimator ($h=0$), $V_0=\Var Y/n=(\E v+\Var g)/n$ and $V_0-V_\star=\Var(g)\big(\frac1n-\frac1{n+N}\big)=\frac{N}{n(n+N)}\Var g$; hence $(V_0-V_{\rm lin})/(V_0-V_\star)=\rho^2$. The width ratio statement follows from $\text{width}\propto\sqrt{V}$. For the expression of $\Gamma$: $V_{\rm lin}/V_\star=1+\Gamma$ with
$\Gamma=\frac{N\Var g\,(1-\rho^2)}{(n+N)\E v+n\Var g}$. \qed

\paragraph{Proof of Corollary~\ref{cor:clip}.}
$V(\lambda f)$ is quadratic in $\lambda$ with leading coefficient $(1/n+1/N)\Var f$ and minimiser $\lambda^\star=\Cov(g,f)/\big((1+n/N)\Var f\big)$, hence $V(\lambda f)=V(\lambda^\star f)+(1/n+1/N)\Var(f)(\lambda-\lambda^\star)^2$. Taking $\lambda=\mathrm{clip}(\lambda^\star,0,1)$ gives the stated excess when $\lambda^\star>1$ (or $<0$). \qed

\paragraph{Remark on efficiency (not a full proof).}
Let $R_k\in\{0,1\}$ indicate labeling, $\pi=n/(n+N)$ known and independent of $(f,Y)$. The observed-data model is then nonparametric for the law of $(f,Y)$, so there is a single influence function, $\varphi=\frac{R}{\pi}(Y-g(f))+g(f)-\theta$, whose variance divided by $n+N$ is $\frac{\E v}{\pi(n+N)}+\frac{\Var g}{n+N}=V_\star$. The estimator $\hat\theta_{cg}$ attains this value, which is why we call $V_\star$ the efficiency bound for this information structure (labels MCAR; $f$ observed everywhere). We use it as the benchmark for the headroom; the claims of the paper do not depend on its optimality beyond the elementary within-class optimality in Proposition~\ref{prop:main}.

\paragraph{Heuristic for the cost of estimating the controls.}
With $p$ estimated control coefficients, the classical regression-adjustment literature gives a relative variance inflation of order $p/n$ for the point estimator, and the in-sample residual variance is biased downwards by a factor of order $(n-p-1)/n$; the plug-in Wald interval is therefore too short by about the same relative amount. We do not derive this here; the data in Section~\ref{sec:exp} are used to check its qualitative implication (undercoverage of the richer control sets at small $n$, partly repaired by the degrees-of-freedom correction).
